\section{Introduction}
Robots that climb or follow staircases, along with assistive devices for
the visually impaired, must reason about three-dimensional structure to
act safely. Legged and wheeled platforms increasingly depend on
egocentric depth sensing to negotiate stairs and other discontinuous
terrain \cite{agarwal2022legged,belter2016adaptive,stelzer2012stereo}.
Shared-autonomy systems likewise demand reliable geometric cues before a
foot or wheel is committed to a surface \cite{marion2017director}. Across
all of these settings the central perceptual question stays the same:
which parts of an observed point cloud form a \emph{tread}, the steppable
horizontal surface on which weight can be safely placed, and which form a
near-vertical \emph{riser} or are simply spurious edge and outlier
returns.

The problem is safety-critical. Labeling a tread as non-steppable only
sacrifices an option, but labeling a riser or edge as steppable invites a
foothold onto a vertical face, which produces slips, stumbles, or falls.
The same hazard motivates stair perception for the visually impaired,
where undetected steps are a recognised fall hazard behind dedicated
stair-detection systems \cite{matsumura2022deep}. Reliable tread
detection therefore needs high recall and, just as importantly, high
precision on the steppable class.

Recent work attacks stair and terrain perception with deep neural
networks that learn point-wise semantics directly from raw clouds
\cite{qi2016pointnet,matsumura2022deep}. These methods are powerful but
data hungry. They require large volumes of densely labeled
three-dimensional scans, which are costly to acquire for the specific
geometry of stairs, and they offer limited interpretability. The tread
sub-problem also lacks controlled, reproducible benchmarks with exact
ground-truth labels, so reported results are hard to compare across
studies. Classical geometric methods are fast, transparent, and
label-free by contrast. Yet despite their maturity in plane segmentation
and traversability estimation
\cite{chavezgarcia2017learning,gomes2023survey}, they remain
under-compared on the staircase setting under matched conditions.

We address this gap with a controlled, fully reproducible study of
classical geometric segmentation for steppable-surface detection on
staircases. Our contributions are as follows:
\begin{itemize}
  \item A parametric, fully labeled staircase point-cloud generator that
        emits exact tread, riser, and other labels while modeling LiDAR
        range jitter and outliers, enabling reproducible evaluation.
  \item A unified three-class problem formulation together with the
        implementation of four classical methods under a common
        orientation-aware pipeline: RANSAC plane fitting, PCA normal
        analysis, DBSCAN with slope confirmation, and a height-histogram
        level detector.
  \item A quantitative comparison of all four methods on the steppable
        (tread) class, accompanied by a controlled noise-sensitivity
        study that characterizes degradation as height noise grows.
  \item A safety-oriented discussion of the results for robotic stair
        traversal, relating precision and intersection-over-union to the
        risk of unsafe footholds.
\end{itemize}

The remainder of this paper is organized as follows: Section~II reviews
related work, Section~III presents the proposed data generator and the
four segmentation methods, Section~IV reports the comparative evaluation
and noise study, and Section~V concludes.
